\begin{lemma}
\label{lemma-limit-is-affine}
Notation and assumptions as in Situation \ref{situation-descent}.
If $X$ is affine, then there exists an $i$ such that $X_i$ is affine.
\end{lemma}

\begin{proof}
Choose $0 \in I$. Choose an affine scheme $U_0$ and a surjective
\'etale morphism $U_0 \to X_0$. Set $U = U_0 \times_{X_0} X$
and $U_i = U_0 \times_{X_0} X_i$ for $i \geq 0$. Since the transition
morphisms are affine, the algebraic spaces $U_i$ and $U$ are affine.
Thus $U \to X$ is an \'etale morphism of affine schemes. Hence we
can write $X = \Spec(A)$, $U = \Spec(B)$ and
$$
B = A[x_1, \ldots, x_n]/(g_1, \ldots, g_n)
$$
such that $\Delta = \det(\partial g_\lambda/\partial x_\mu)$ is invertible
in $B$, see Algebra, Lemma \ref{algebra-lemma-etale-standard-smooth}.
Set $A_i = \mathcal{O}_{X_i}(X_i)$. We have $A = \colim A_i$ by
Lemma \ref{lemma-descend-section}. After increasing $0$ we may assume
we have $g_{1, i}, \ldots, g_{n, i} \in A_i[x_1, \ldots, x_n]$ mapping to
$g_1, \ldots, g_n$. Set
$$
B_i = A_i[x_1, \ldots, x_n]/(g_{1, i}, \ldots, g_{n, i})
$$
for all $i \geq 0$. Increasing $0$ if necessary we may assume that
$\Delta_i = \det(\partial g_{\lambda, i}/\partial x_\mu)$ is invertible
in $B_i$ for all $i \geq 0$. Thus $A_i \to B_i$ is an \'etale ring map.
After increasing $0$ we may assume also that
$\Spec(B_i) \to \Spec(A_i)$ is surjective, see
Limits, Lemma \ref{limits-lemma-descend-surjective}. Increasing
$0$ yet again we may choose elements
$h_{1, i}, \ldots, h_{n, i} \in \mathcal{O}_{U_i}(U_i)$ which map to the
classes of $x_1, \ldots, x_n$ in $B = \mathcal{O}_U(U)$ and such that
$g_{\lambda, i}(h_{\nu, i}) = 0$ in $\mathcal{O}_{U_i}(U_i)$. Thus
we obtain a commutative diagram
\begin{equation}
\label{equation-to-show-cartesian}
\vcenter{
\xymatrix{
X_i \ar[d] & U_i \ar[l] \ar[d] \\
\Spec(A_i) & \Spec(B_i) \ar[l]
}
}
\end{equation}
By construction $B_i = B_0 \otimes_{A_0} A_i$ and
$B = B_0 \otimes_{A_0} A$. Consider the morphism
$$
f_0 : U_0 \longrightarrow X_0 \times_{\Spec(A_0)} \Spec(B_0)
$$
This is a morphism of quasi-compact and quasi-separated algebraic spaces
representable, separated and \'etale over $X_0$. The base change of $f_0$
to $X$ is an isomorphism by our choices. Hence
Lemma \ref{lemma-descend-equality}
guarantees that there exists an $i$ such that the base change of $f_0$
to $X_i$ is an isomorphism, in other words the diagram
(\ref{equation-to-show-cartesian}) is cartesian. Thus
Descent, Lemma \ref{descent-lemma-descent-data-sheaves}
applied to the fppf covering $\{\Spec(B_i) \to \Spec(A_i)\}$
combined with Descent, Lemma \ref{descent-lemma-affine}
give that $X_i \to \Spec(A_i)$ is representable by a scheme
affine over $\Spec(A_i)$ as desired. (Of course it then also follows
that $X_i = \Spec(A_i)$ but we don't need this.)
\end{proof}
